\section{Relation to existing theory}
\label{sec:comparison}

The comparison has two parts.  At exponent one, RL recovers familiar
quadratic operator inequalities and their PPA distance estimates.  At an
exponent below one, the square-root example separates the present assumptions
from specific quadratic certificates on the same neighborhood.  All of the
operator frameworks compared below allow set-valued operators.

\subsection{The quadratic endpoint}

\begin{proposition}[Identification of the exponent-one condition]
\label{prop:linear-identifications}
Fix a graph block and $\lambda>0$.  For each tested pair of its graph points,
write $a=u-v$ and $b=u^*-v^*$.  The following statements hold on the same
pair range as the RL condition.
\begin{enumerate}
\item The condition $\mathrm{RL}(1,L)$ is equivalent to
\begin{equation}
 \lambda\langle a,b\rangle
 \ge \frac{1-L^2}{2(1+L^2)}
       \bigl(\|a\|^2+\lambda^2\|b\|^2\bigr).
 \label{eq:linear-rl-balanced}
\end{equation}
Consequently $\mathrm{RL}(1,1)$ is precisely monotonicity on the block.
In the notation of $(\mu,\nu)$-semimonotonicity,
\begin{equation}
 \mu=\frac{1-L^2}{2\lambda(1+L^2)},
 \qquad \nu=\lambda^2\mu.
 \label{eq:balanced-slice-parameters}
\end{equation}
Thus $\mathrm{RL}(1,L)$ is the balanced parameter slice of that inequality.

\item The scaled submonotonicity condition
\[
 \lambda\langle a,b\rangle
 \ge-\tau\|a+\lambda b\|^2,\qquad \tau\ge0,
\]
is equivalent to $\mathrm{RL}(1,L)$ with
\begin{equation}
 L^2=1+4\tau.
 \label{eq:luke-tam-normalization}
\end{equation}

\item If the block is $h$-hypomonotone, namely
\[
 \langle a,b\rangle\ge-h\|a\|^2,
 \qquad h\ge0,
\]
and $\lambda h<1$, then it satisfies $\mathrm{RL}(1,L_h)$ with
\begin{equation}
 L_h=\frac{1+\lambda h}{1-\lambda h}.
 \label{eq:hypo-to-rl}
\end{equation}
\end{enumerate}
\end{proposition}

\begin{proof}
At $\gamma=1$, the RL inner-product inequality is
\[
 4\lambda\langle a,b\rangle
 \ge(1-L^2)\|a+\lambda b\|^2.
\]
Expanding the squared norm and collecting the inner-product terms proves
\eqref{eq:linear-rl-balanced}.  Setting $L=1$ proves the monotone endpoint;
dividing by $\lambda$ gives \eqref{eq:balanced-slice-parameters}.  Substituting
$1-L^2=-4\tau$ proves the second assertion.

For the last assertion, set $c=\lambda b$ and $\eta=\lambda h<1$.
Hypomonotonicity and Cauchy--Schwarz give
\[
 (1-\eta)\|a\|^2
 \le\langle a,a+c\rangle
 \le\|a\|\,\|a+c\|,
\]
so $\|a\|\le\|a+c\|/(1-\eta)$.  Therefore
\[
 \begin{split}
 \|a\|^2+\|c\|^2
 &=\|a+c\|^2-2\langle a,c\rangle\\
 &\le\left(1+\frac{2\eta}{(1-\eta)^2}\right)\|a+c\|^2
 =\frac{1+L_h^2}{2}\|a+c\|^2.
 \end{split}
\]
This is the all-pairs $J$-side inequality with exponent one, under the
identification $x=u+\lambda u^*$, $y=v+\lambda v^*$.
\end{proof}

The semimonotonicity parameters in
\eqref{eq:balanced-slice-parameters} use the convention of
\cite{EvensPasLatafatPatrinos2025SemimonotonePPA};
\eqref{eq:luke-tam-normalization} gives the scaled Luke--Tam condition
\cite{LukeTam2025GeneralizedMonotonicity}.  Our local range condition (A1)
uses the same coverage mechanism as
\cite[Assumption~2(d)(ii)]{LukeTam2025GeneralizedMonotonicity} as its
well-posedness component.  The extension occurs in (A2): the exponent-one
quadratic graph inequality is replaced by the nonlinear RL order, while
(A3)--(A4) match that order with the output error bound.

\subsection{Error bounds and the linear convergence factor}

The regularity assumption concerns the original operator at the proximal
output:
\[
 \dist(y,S)\le\psi\bigl(\dist(0,F(y))\bigr).
\]
It can be verified independently of RL.  In particular, an available
metric $q$-subregularity estimate takes the form
\begin{equation}
 \dist(y,S)\le K\dist(0,F(y))^q.
 \label{eq:power-eb-comparison}
\end{equation}
The orders $q>0$, including $q>1$, belong to the established subregularity
theory \cite{LiMordukhovich2012HolderSubregularityPPA,
MordukhovichOuyang2015HigherOrderSubregularity}.  Substituting
$\psi(t)=Kt^q$ into the direct recursion gives the coefficient
$K(L/(2\lambda))^q$ and the leading power $d^{\gamma q}$ for
$0<\gamma<1$, $L>0$.  This explains exactly which independently obtained
error bound can be used: $\gamma q>1$ gives local compatibility, while
$\gamma q=1$ requires $K(L/(2\lambda))^q<1$.

General gauges also enter the abstract framework of
\cite[Theorems~4.3--4.4]{LiMordukhovichZhu2026GeneralizedSubregularity}.
That framework concerns subgradient sequences with an objective descent
condition, a surrogate for successive change, a relative subgradient-error
bound, a continuity condition, and an accumulation point.  Its regularity is
written as a gauge applied to distance and is combined with the surrogate and
relative-error moduli to obtain the rate map.  Our theorem does not assume an
underlying objective, descent, or a pre-existing accumulation point.  Instead,
the RL inequality supplies the step estimate, the proximal inclusion supplies
the output residual, and an initial-radius budget proves localization and the
existence of a limit.  Thus the two frameworks overlap for suitable
subgradient PPA specializations, but neither set of assumptions follows from
the other at the present level of generality.  The gauge assumption used here
is an imported regularity condition; the square-root example verifies it by
direct residual calculation.

Corollary~\ref{cor:linear-slice} gives the coupled factor
\[
 \kappa_E=
 \sqrt{\frac{1+L^2}{2}}\,
 \frac{\rho}{\sqrt{\rho^2+\lambda^2}}.
\]
At $L=1$, this is
$\rho/\sqrt{\rho^2+\lambda^2}$, the distance factor in
Leventhal's Theorem~3.1 \cite{Leventhal2009MetricSubregularityPPA}.
Under the Luke--Tam normalization $L^2=1+4\tau$,
\begin{equation}
 \kappa_{\rm coupled}^2
 =(1+2\tau)\frac{\rho^2}{\rho^2+\lambda^2},
 \qquad
 \kappa_{\rm coupled}<1
 \quad\Longleftrightarrow\quad
 2\tau\rho^2<\lambda^2.
 \label{eq:coupled-linear-factor}
\end{equation}
For the common scaled submonotonicity inequality, the estimate displayed in
\cite{LukeTam2025GeneralizedMonotonicity} is
\[
 \kappa_{\rm LT}^2
 =1+2\tau-\left(\frac{\lambda}{\lambda+\rho}\right)^2,
 \qquad
 \tau(1+\rho/\lambda)^2<\frac12.
\]
Keeping the error bound at the output gives
\eqref{eq:coupled-linear-factor}, with a smaller factor and a weaker scalar
restriction.  This comparison is confined to the common transition
inequalities; coverage and localization must also hold in either local
convergence theorem.

\subsection{What the multivalued example separates}

Almost-averaged mapping theory explicitly allows set-valued iteration maps
\cite{LukeThaoTam2018Quantitative}.  Generalized-monotonicity PPA theory also
allows set-valued operators
\cite{LukeTam2025GeneralizedMonotonicity,
EvensPasLatafatPatrinos2025SemimonotonePPA}.  The distinction in the following
proposition is the order of the estimate near the branch junction.

\begin{proposition}[Failure of specified quadratic certificates]
\label{prop:seam-quadratic-separation}
Let $F$ be the square-root operator in
Theorem~\ref{thm:square-root-seam}, let $J=J_F$, and fix
$\bar z=(\bar p,0)\in S$.  On every full neighborhood of $\bar z$:
\begin{enumerate}
\item there is no finite $C$ such that
\begin{equation}
 \|Jz-Jw\|^2+
 \|(z-Jz)-(w-Jw)\|^2\le C\|z-w\|^2
 \label{eq:excluded-quadratic-J}
\end{equation}
for all inputs $z,w$ in that neighborhood;
\item there is no finite $C$ such that
$\|Jz-\bar z\|\le C\|z-\bar z\|$ for all nearby inputs;
\item there is no finite local hypomonotonicity modulus for the graph of $F$
near $(\bar z,0)$;
\item there is no fixed finite matrix $V$ such that
\begin{equation}
 \langle u-\bar z,u^*\rangle\ge\langle u^*,Vu^*\rangle
 \label{eq:excluded-fixed-anchor-matrix}
\end{equation}
for all graph points $(u,u^*)$ sufficiently near $(\bar z,0)$.
\end{enumerate}
\end{proposition}

\begin{proof}
The resolvent is
$J(p,q)=(p+\sqrt{|q|},|q|/4)$.  Take
$z_t=(\bar p,t)$ and $w_t=(\bar p,-4t)$, $t>0$.  Direct calculation gives
\[
 \|z_t-w_t\|^2=25t^2,
 \qquad
 \|Jz_t-Jw_t\|^2+
 \|(z_t-Jz_t)-(w_t-Jw_t)\|^2
 =2t+\frac{269}{8}t^2.
\]
Their ratio diverges, proving the first assertion.  For the second,
\[
 \frac{\|Jz_t-\bar z\|}{\|z_t-\bar z\|}
 =\frac{\sqrt{t+t^2/16}}{t}\longrightarrow+\infty.
\]
To prove the third assertion, use the graph points
$u_t=(\bar p+t,t)$ and $u_t^*=(-2\sqrt t,3t)$, together with
$(\bar z,0)$.  They satisfy
\[
 \frac{\langle u_t-\bar z,u_t^*\rangle}
      {\|u_t-\bar z\|^2}
 =-\frac1{\sqrt t}+\frac32\longrightarrow-\infty.
\]
Finally, replace $u_t$ by $(\bar p+t^{1/4},t)$, keeping the same $u_t^*$.
Then
\[
 \frac{\langle u_t-\bar z,u_t^*\rangle}{\|u_t^*\|^2}
 =\frac{-2t^{3/4}+3t^2}{4t+9t^2}\longrightarrow-\infty.
\]
For a fixed finite matrix $V$, the quotient
$\langle u_t^*,Vu_t^*\rangle/\|u_t^*\|^2$ is bounded below by
$-\|V\|$.  This proves the last assertion.
\end{proof}

The four assertions exclude the displayed quadratic certificates on the same
full neighborhood.  They do not compare the entire semimonotone, weak-Minty,
almost-averaged, or partial-regularity frameworks, which may use different
relative domains or degenerate tests.  In contrast, the H\"older $J$-side
estimate and the quadratic output error bound of
Theorem~\ref{thm:square-root-seam} satisfy the RLEB compatibility condition.
